% =============================================================================
% SECTION <7>: TECHNICAL REQUIREMENTS
% =============================================================================
\subsection{7. Technical Requirements}

This section establishes the normative, verifiable Level-2 technical requirements governing Configuration Item \textbf{\texttt{CI-03}: Space Segment (\abbr{SSEG}, \abbr{WBS} Element \texttt{123456.03})} and its integrated orbital constellation comprising seventy-two (72) Low Earth Orbit tracking satellites (\textbf{\texttt{SATL}, \abbr{WBS} \texttt{123456.03.01}}) at 1,000~km circular altitude across 8 orbital planes and six (6) Geostationary Earth Orbit detection satellites (\textbf{\texttt{SATG}, \abbr{WBS} \texttt{123456.03.02}}) at 35,786~km altitude (\texttt{R-005}). Segment-level design definitions and technical budgets are documented in \texttt{EWAS-HSB-SSEG-DDF-03} (\textbf{RD-01}) and \texttt{EWAS-HSB-SSEG-DDJF-03} (\textbf{RD-02}), while detailed internal spacecraft and subsystem requirements are allocated to the subordinate Level-3 Technical Requirements Specifications for \texttt{SATL} (\texttt{EWAS-HSB-SATL-TRS-03.01}, \textbf{RD-03}) and \texttt{SATG} (\texttt{EWAS-HSB-SATG-TRS-03.02}, \textbf{RD-04}).

Every technical requirement specified in Subsections 7.1 through 7.8 is uniquely identified by a persistent code \texttt{SSEG-REQ-<CAT>-<NNN>} and formulated using the mandatory verb \textbf{\texttt{shall}}. Each requirement traces backward to the parent system requirements (\texttt{R-001} through \texttt{R-049}) of \textbf{AD-01} and is assigned one or more formal verification methods defined in the Verification Plan (\textbf{AD-02}): \textbf{\texttt{T}} (\emph{Test}), \textbf{\texttt{A}} (\emph{Analysis}), \textbf{\texttt{R}} (\emph{Review of Design}), and \textbf{\texttt{I}} (\emph{Inspection}).

% -----------------------------------------------------------------------------
% 7.1 FUNCTIONAL AND MISSION SCOPE REQUIREMENTS
% -----------------------------------------------------------------------------
\subsubsection{7.1 Functional and Mission Scope Requirements (\texttt{SSEG-REQ-FUN})}

Table~\ref{tab:req_fun} defines the top-level functional, non-offensive civil-protection boundary, and warning-chain execution requirements allocated to the \abbr{EWAS} Space Segment.

{\footnotesize\setlength{\tabcolsep}{4pt}\renewcommand{\arraystretch}{1.15}%
\arrayrulecolor{TableBorderBlue}%
\rowcolors{2}{white}{TableRowLightBlue}%
\setlength{\LTcapwidth}{\textwidth}%
\begin{longtable}{|>{\raggedright\arraybackslash}p{0.18\textwidth}|>{\raggedright\arraybackslash}p{0.58\textwidth}|>{\raggedright\arraybackslash}p{0.12\textwidth}|>{\raggedright\arraybackslash}p{0.06\textwidth}|}
\caption{Functional and Mission Scope Technical Requirements (\texttt{SSEG-REQ-FUN})}\label{tab:req_fun}\\
\hline
\rowcolor{TableHeaderBlue}
\thead{Req.\ ID \& Title} & \thead{Normative Requirement Statement (\texttt{shall})} & \thead{Parent Req.} & \thead{Verif.} \\ \hline
\endfirsthead
\rowcolor{white}\multicolumn{4}{l}{\footnotesize\textit{Table \thetable\ (continued): Functional and Mission Scope Technical Requirements (\texttt{SSEG-REQ-FUN})}} \\[0.3ex]
\hline
\rowcolor{TableHeaderBlue}
\thead{Req.\ ID \& Title} & \thead{Normative Requirement Statement (\texttt{shall})} & \thead{Parent Req.} & \thead{Verif.} \\ \hline
\endhead
\rowcolor{white}\multicolumn{4}{r}{\footnotesize\textit{Continued on next page}} \\
\endfoot
\endlastfoot
\textbf{\texttt{SSEG-REQ-FUN-001}} \newline \emph{Non-Offensive Civil-Protection Operation} & The \abbr{SSEG} hardware, flight software, and downlinked data products \textbf{shall} exclusively support ballistic missile launch detection, stereoscopic trajectory tracking, and impact-area warning generation for civil protection, and \textbf{shall} exclude any weapon-guidance, fire-control, or kinetic-interceptor command capabilities. & \texttt{R-001}, \newline \texttt{R-048} & \texttt{R} \\ \hline
\textbf{\texttt{SSEG-REQ-FUN-002}} \newline \emph{Scope Exclusion of Yield and Damage Assessment} & The \abbr{SSEG} \textbf{shall} not be required to perform warhead nuclear or conventional yield estimation, definitive missile-variant identification, or post-impact terrestrial damage assessment. & \texttt{R-002} & \texttt{R} \\ \hline
\textbf{\texttt{SSEG-REQ-FUN-003}} \newline \emph{Dual-Tier Detection and Stereoscopic Tracking Chain} & The \abbr{SSEG} \textbf{shall} execute continuous boost-phase infrared detection via the \abbr{GEO} layer (\texttt{SATG}), route tracking commands via the mandatory \abbr{OISL} network, execute cued dual-satellite stereoscopic tracking via the \abbr{LEO} layer (\texttt{SATL}), compute an onboard preliminary 3D trajectory estimate, and downlink both timestamped \abbr{LOS} observation packages and the preliminary trajectory estimate to the Ground Segment (\abbr{WBS} \texttt{123456.02}). & \texttt{R-007}, \texttt{R-009}, \texttt{R-023}, \texttt{R-042}, \texttt{R-044} & \texttt{T, A} \\ \hline
\textbf{\texttt{SSEG-REQ-FUN-004}} \newline \emph{Baltic Region and Extra-Regional Threat Surveillance} & The \abbr{SSEG} \textbf{shall} monitor both the primary Baltic Sea operational theatre and extra-regional launch corridors across adjacent and distant theatres capable of launching ballistic threats toward the Baltic region. & \texttt{R-006}, \newline \texttt{R-049} & \texttt{A} \\ \hline
\textbf{\texttt{SSEG-REQ-FUN-005}} \newline \emph{Onboard Proton Despiking and Plume Discrimination} & Both \texttt{SATG} and \texttt{SATL} \textbf{shall} perform real-time onboard proton despiking, background clutter suppression, and dual-band spectral/kinematic discrimination to distinguish rocket exhaust plumes from solar specular glints, terrestrial fires, aurorae, and radiation-induced detector transients. & \texttt{R-019}, \newline \texttt{R-022} & \texttt{T, A} \\ \hline
\textbf{\texttt{SSEG-REQ-FUN-006}} \newline \emph{Archival Telemetry and Provenance Formatting} & All downlinked \abbr{SSEG} event reports, \abbr{LOS} observation packages, preliminary trajectory estimates, and housekeeping records \textbf{shall} include GNSS-synchronized UTC timestamps, spacecraft ephemeris and attitude covariance metadata, and data-quality flags formatted for at least 5-year sovereign archival in the Ground Segment Data Archiving System (\texttt{123456.02.01}) exclusively for calibration, training, verification, and algorithm refinement. & \texttt{R-004}, \texttt{R-036}, \texttt{R-047} & \texttt{R, T} \\ \hline
\end{longtable}}

% -----------------------------------------------------------------------------
% 7.2 CONSTELLATION ARCHITECTURE REQUIREMENTS
% -----------------------------------------------------------------------------
\subsubsection{7.2 Constellation Architecture Requirements (\texttt{SSEG-REQ-MIS})}

Table~\ref{tab:req_mis} specifies the Level-2 orbital configuration, continuous surveillance coverage, and stereoscopic viewing geometry requirements for the hybrid 6~\abbr{GEO} + 72~\abbr{LEO} constellation (\abbr{WBS} \texttt{123456.03.01.07} and \texttt{123456.03.02.07}).

{\footnotesize\setlength{\tabcolsep}{4pt}\renewcommand{\arraystretch}{1.15}%
\arrayrulecolor{TableBorderBlue}%
\rowcolors{2}{white}{TableRowLightBlue}%
\setlength{\LTcapwidth}{\textwidth}%
\begin{longtable}{|>{\raggedright\arraybackslash}p{0.18\textwidth}|>{\raggedright\arraybackslash}p{0.58\textwidth}|>{\raggedright\arraybackslash}p{0.12\textwidth}|>{\raggedright\arraybackslash}p{0.06\textwidth}|}
\caption{Constellation Architecture Requirements (\texttt{SSEG-REQ-MIS})}\label{tab:req_mis}\\
\hline
\rowcolor{TableHeaderBlue}
\thead{Req.\ ID \& Title} & \thead{Normative Requirement Statement (\texttt{shall})} & \thead{Parent Req.} & \thead{Verif.} \\ \hline
\endfirsthead
\rowcolor{white}\multicolumn{4}{l}{\footnotesize\textit{Table \thetable\ (continued): Constellation Architecture Requirements (\texttt{SSEG-REQ-MIS})}} \\[0.3ex]
\hline
\rowcolor{TableHeaderBlue}
\thead{Req.\ ID \& Title} & \thead{Normative Requirement Statement (\texttt{shall})} & \thead{Parent Req.} & \thead{Verif.} \\ \hline
\endhead
\rowcolor{white}\multicolumn{4}{r}{\footnotesize\textit{Continued on next page}} \\
\endfoot
\endlastfoot
\textbf{\texttt{SSEG-REQ-MIS-001}} \newline \emph{Hybrid Constellation Architecture} & The \abbr{SSEG} \textbf{shall} consist of six (6) active \abbr{GEO} detection satellites (\texttt{SATG}) operating at 35,786~km altitude and seventy-two (72) active \abbr{LEO} tracking satellites (\texttt{SATL}) operating at 1,000~km circular altitude across 8 orbital planes. & \texttt{R-005} & \texttt{A, R} \\ \hline
\textbf{\texttt{SSEG-REQ-MIS-002}} \newline \emph{Continuous 24/7 Surveillance Coverage} & The \abbr{SSEG} constellation \textbf{shall} provide continuous, uninterrupted 24-hour spatial and temporal surveillance of the designated Baltic monitoring region and relevant extra-regional launch corridors without operational blind windows. & \texttt{R-006}, \texttt{R-008}, \texttt{R-049} & \texttt{A} \\ \hline
\textbf{\texttt{SSEG-REQ-MIS-003}} \newline \emph{Boost-Phase Detection Coverage} & The \texttt{SATG} geostationary layer, supplemented at high northern latitudes by the \texttt{SATL} layer, \textbf{shall} maintain continuous observation geometry over potential launch zones to detect ballistic missiles during their boost phase. & \texttt{R-007} & \texttt{A} \\ \hline
\textbf{\texttt{SSEG-REQ-MIS-004}} \newline \emph{Multi-Satellite Stereoscopic Tracking Geometry} & The 72-satellite, 8-plane \texttt{SATL} constellation at 1,000~km \textbf{shall} be phased (in accordance with the astrodynamic definition in \texttt{EWAS-HSB-SSEG-DDF-03}, \textbf{RD-01}) such that at least two (2) \texttt{SATL} spacecraft simultaneously maintain mutually intersecting lines of sight to any boost- or post-boost-phase missile threat affecting the monitored theatre to enable 3D stereoscopic triangulation. & \texttt{R-007}, \texttt{R-008}, \texttt{R-023} & \texttt{A} \\ \hline
\textbf{\texttt{SSEG-REQ-MIS-005}} \newline \emph{Constellation Maintenance Deconfliction} & Orbit maintenance and station-keeping manoeuvres across the constellation \textbf{shall} be scheduled and staggered such that continuous surveillance and stereoscopic tracking availability are never degraded. & \texttt{R-008}, \newline \texttt{R-015} & \texttt{A, R} \\ \hline
\end{longtable}}

% -----------------------------------------------------------------------------
% 7.3 SPACE SEGMENT PERFORMANCE REQUIREMENTS
% -----------------------------------------------------------------------------
\subsubsection{7.3 Space Segment Performance Requirements (\texttt{SSEG-REQ-PERF})}

Table~\ref{tab:req_perf} establishes the Space Segment launch detection probability, end-to-end warning latency contribution (within the $\le 60\text{ s}$ system mandate of \texttt{R-011}), false-alarm suppression, and operational availability requirements.

{\footnotesize\setlength{\tabcolsep}{4pt}\renewcommand{\arraystretch}{1.15}%
\arrayrulecolor{TableBorderBlue}%
\rowcolors{2}{white}{TableRowLightBlue}%
\setlength{\LTcapwidth}{\textwidth}%
\begin{longtable}{|>{\raggedright\arraybackslash}p{0.18\textwidth}|>{\raggedright\arraybackslash}p{0.58\textwidth}|>{\raggedright\arraybackslash}p{0.12\textwidth}|>{\raggedright\arraybackslash}p{0.06\textwidth}|}
\caption{Space Segment Performance Requirements (\texttt{SSEG-REQ-PERF})}\label{tab:req_perf}\\
\hline
\rowcolor{TableHeaderBlue}
\thead{Req.\ ID \& Title} & \thead{Normative Requirement Statement (\texttt{shall})} & \thead{Parent Req.} & \thead{Verif.} \\ \hline
\endfirsthead
\rowcolor{white}\multicolumn{4}{l}{\footnotesize\textit{Table \thetable\ (continued): Space Segment Performance Requirements (\texttt{SSEG-REQ-PERF})}} \\[0.3ex]
\hline
\rowcolor{TableHeaderBlue}
\thead{Req.\ ID \& Title} & \thead{Normative Requirement Statement (\texttt{shall})} & \thead{Parent Req.} & \thead{Verif.} \\ \hline
\endhead
\rowcolor{white}\multicolumn{4}{r}{\footnotesize\textit{Continued on next page}} \\
\endfoot
\endlastfoot
\textbf{\texttt{SSEG-REQ-PERF-001}} \newline \emph{Boost-Phase Launch Detection Probability} & The \abbr{SSEG} \textbf{shall} achieve a boost-phase missile launch probability of detection exceeding $95\%$ ($P_d > 95\%$) for missile threats defined within the mission scope. & \texttt{R-009}, \newline \texttt{R-010} & \texttt{T, A} \\ \hline
\textbf{\texttt{SSEG-REQ-PERF-002}} \newline \emph{End-to-End Warning Latency Contribution} & \abbr{SSEG} boost-phase detection, candidate event reporting, inter-satellite tracking command routing, cued \abbr{LEO} stereoscopic tracking, onboard preliminary 3D trajectory estimation, and observation downlink \textbf{shall} execute within the Space Segment latency budget allocated in \texttt{EWAS-HSB-SSEG-DDJF-03} (\textbf{RD-02}) such that the total elapsed time from first \abbr{GEO} detection to delivery of the preliminary warning product at the authorized national authority interface does not exceed 60~seconds. & \texttt{R-011}, \texttt{R-033}, \texttt{R-042}, \texttt{R-044}, \texttt{R-045} & \texttt{T, A} \\ \hline
\textbf{\texttt{SSEG-REQ-PERF-003}} \newline \emph{Onboard False-Alarm Suppression} & \abbr{SSEG} onboard proton despiking and background clutter discrimination \textbf{shall} suppress non-threat infrared sources (including solar glints, cloud reflections, terrestrial fires, and radiation transients) in accordance with the false-alarm budgets defined in \textbf{RD-02} while maintaining $P_d > 95\%$. & \texttt{R-010}, \texttt{R-019}, \texttt{R-022} & \texttt{T, A} \\ \hline
\textbf{\texttt{SSEG-REQ-PERF-004}} \newline \emph{Space Segment Operational Availability} & The integrated \abbr{SSEG} constellation \textbf{shall} achieve an operational availability supporting an overall \abbr{EWAS} system availability exceeding $99.9\%$. & \texttt{R-015}, \newline \texttt{R-017} & \texttt{A} \\ \hline
\end{longtable}}

% -----------------------------------------------------------------------------
% 7.4 SUBORDINATE FLIGHT SYSTEM ALLOCATION REQUIREMENTS
% -----------------------------------------------------------------------------
\subsubsection{7.4 Subordinate Flight System Allocation Requirements (\texttt{SSEG-REQ-ALLOC})}

In accordance with the EWAS specification tree (Table~\ref{tab:sseg_wbs_scope}), this subsection defines the top-level functional, pointing, processing, and technical-budget compliance requirements allocated by the Space Segment (\texttt{CI-03}) to its subordinate Level-3 flight system Configuration Items: Satellite \abbr{LEO} (\texttt{SATL}, \abbr{WBS} \texttt{123456.03.01}) and Satellite \abbr{GEO} (\texttt{SATG}, \abbr{WBS} \texttt{123456.03.02}). Detailed spacecraft-level and subsystem-level numerical requirements (governing internal Payload optics/cryogenics, \abbr{ADCS} sensors/actuators, \abbr{OBDH} hardware/memory, Structures \& Mechanisms, Thermal Control, Electrical Power, and Propulsion subsystems) are specified in \texttt{EWAS-HSB-SATL-TRS-03.01} (\textbf{RD-03}) and \texttt{EWAS-HSB-SATG-TRS-03.02} (\textbf{RD-04}).

{\footnotesize\setlength{\tabcolsep}{4pt}\renewcommand{\arraystretch}{1.15}%
\arrayrulecolor{TableBorderBlue}%
\rowcolors{2}{white}{TableRowLightBlue}%
\setlength{\LTcapwidth}{\textwidth}%
\begin{longtable}{|>{\raggedright\arraybackslash}p{0.18\textwidth}|>{\raggedright\arraybackslash}p{0.58\textwidth}|>{\raggedright\arraybackslash}p{0.12\textwidth}|>{\raggedright\arraybackslash}p{0.06\textwidth}|}
\caption{Subordinate Flight System Allocation Requirements (\texttt{SSEG-REQ-ALLOC})}\label{tab:req_alloc}\\
\hline
\rowcolor{TableHeaderBlue}
\thead{Req.\ ID \& Title} & \thead{Normative Requirement Statement (\texttt{shall})} & \thead{Parent Req.} & \thead{Verif.} \\ \hline
\endfirsthead
\rowcolor{white}\multicolumn{4}{l}{\footnotesize\textit{Table \thetable\ (continued): Subordinate Flight System Allocation Requirements (\texttt{SSEG-REQ-ALLOC})}} \\[0.3ex]
\hline
\rowcolor{TableHeaderBlue}
\thead{Req.\ ID \& Title} & \thead{Normative Requirement Statement (\texttt{shall})} & \thead{Parent Req.} & \thead{Verif.} \\ \hline
\endhead
\rowcolor{white}\multicolumn{4}{r}{\footnotesize\textit{Continued on next page}} \\
\endfoot
\endlastfoot
\textbf{\texttt{SSEG-REQ-ALLOC-001}} \newline \emph{Infrared Sensing Payload Allocation} & Each \texttt{SATG} and \texttt{SATL} flight system \textbf{shall} carry dual-band infrared sensing payloads for detecting and tracking missile exhaust plume signatures against background environmental conditions in compliance with the optical and radiometric specifications allocated in \texttt{EWAS-HSB-SATL-TRS-03.01} (\textbf{RD-03}) and \texttt{EWAS-HSB-SATG-TRS-03.02} (\textbf{RD-04}). & \texttt{R-007}, \texttt{R-009}, \texttt{R-020}, \texttt{R-021} & \texttt{T, A} \\ \hline
\textbf{\texttt{SSEG-REQ-ALLOC-002}} \newline \emph{Independent Optical Payload LOS Lock-On and Pointing Stability Allocation} & Upon receipt of a tracking cue, each cued \texttt{SATL} flight system \textbf{shall} independently steer its payload optical Line of Sight (\abbr{LOS}) to achieve target lock-on within $< 5\text{ seconds}$ and maintain optical \abbr{LOS} pointing stability strictly below $< 2\text{ arcsec/s}$ without slewing the spacecraft bus or disturbing inter-satellite and ground communication links, while maintaining the attitude, ephemeris, and time-synchronization accuracy specified in \textbf{RD-03}. & \texttt{R-012}, \texttt{R-023}, \texttt{R-024} & \texttt{T, A} \\ \hline
\textbf{\texttt{SSEG-REQ-ALLOC-003}} \newline \emph{Onboard Event Filtering and Stereoscopic 3D Trajectory Computation Allocation} & Each flight system \textbf{shall} perform real-time onboard proton despiking, background filtering, and preliminary event characterization (\texttt{R-022}), and during cued tracking the two selected \texttt{SATL} flight systems \textbf{shall} exchange timestamped \abbr{LOS} observation vectors via the mandatory \abbr{OISL} network to compute an \emph{Onboard Preliminary Trajectory} estimate and downlink both the timestamped \emph{LEO Observation Packages} (with covariance metadata) and the preliminary trajectory estimate to the Ground Segment. & \texttt{R-012}, \texttt{R-019}, \texttt{R-022}, \texttt{R-023}, \texttt{R-042}, \texttt{R-045} & \texttt{T, A} \\ \hline
\textbf{\texttt{SSEG-REQ-ALLOC-004}} \newline \emph{Stray-Light Mitigation and Contamination Clearance Allocation} & Both \texttt{SATL} and \texttt{SATG} flight systems \textbf{shall} incorporate optical stray-light mitigation baffles and thruster/outgassing clearance geometries as specified in \textbf{RD-03} and \textbf{RD-04} to prevent stray radiation or molecular contamination from degrading sensor performance. & \texttt{R-020}, \newline \texttt{R-025} & \texttt{T, A} \\ \hline
\textbf{\texttt{SSEG-REQ-ALLOC-005}} \newline \emph{Flight System Technical Budget and Dispenser Envelope Allocation} & Each \texttt{SATL} and \texttt{SATG} flight system \textbf{shall} conform to the mass, electrical power, thermal, structural stiffness, and propulsion $\Delta v$ technical budgets established in \texttt{EWAS-HSB-SSEG-DDF-03} (\textbf{RD-01}) and \texttt{EWAS-HSB-SSEG-DDJF-03} (\textbf{RD-02}), ensuring compatibility with multi-satellite plane-batch dispenser deployment (\texttt{SATL}) and geostationary launch vehicles (\texttt{SATG}). & \texttt{R-005}, \texttt{R-015}, \texttt{R-016}, \texttt{R-018} & \texttt{T, A} \\ \hline
\end{longtable}}

% -----------------------------------------------------------------------------
% 7.5 SPACE SEGMENT COMMUNICATIONS AND MANDATORY OISL NETWORK REQUIREMENTS
% -----------------------------------------------------------------------------
\subsubsection{7.5 Space Segment Communications and Mandatory OISL Network Requirements (\texttt{SSEG-REQ-ISL})}

Table~\ref{tab:req_isl} defines the segment-level requirements for the mandatory Optical Inter-Satellite Link (\abbr{OISL}) network, Space-to-Ground \abbr{RF} \abbr{TT\&C} and mission links (\abbr{WBS} \texttt{123456.03.01.06.02} and \texttt{123456.03.02.06.02}), cryptographic security, and electronic warfare resilience.

{\footnotesize\setlength{\tabcolsep}{4pt}\renewcommand{\arraystretch}{1.15}%
\arrayrulecolor{TableBorderBlue}%
\rowcolors{2}{white}{TableRowLightBlue}%
\setlength{\LTcapwidth}{\textwidth}%
\begin{longtable}{|>{\raggedright\arraybackslash}p{0.18\textwidth}|>{\raggedright\arraybackslash}p{0.58\textwidth}|>{\raggedright\arraybackslash}p{0.12\textwidth}|>{\raggedright\arraybackslash}p{0.06\textwidth}|}
\caption{Communications and Mandatory OISL Network Requirements (\texttt{SSEG-REQ-ISL})}\label{tab:req_isl}\\
\hline
\rowcolor{TableHeaderBlue}
\thead{Req.\ ID \& Title} & \thead{Normative Requirement Statement (\texttt{shall})} & \thead{Parent Req.} & \thead{Verif.} \\ \hline
\endfirsthead
\rowcolor{white}\multicolumn{4}{l}{\footnotesize\textit{Table \thetable\ (continued): Communications and Mandatory OISL Network Requirements (\texttt{SSEG-REQ-ISL})}} \\[0.3ex]
\hline
\rowcolor{TableHeaderBlue}
\thead{Req.\ ID \& Title} & \thead{Normative Requirement Statement (\texttt{shall})} & \thead{Parent Req.} & \thead{Verif.} \\ \hline
\endhead
\rowcolor{white}\multicolumn{4}{r}{\footnotesize\textit{Continued on next page}} \\
\endfoot
\endlastfoot
\textbf{\texttt{SSEG-REQ-ISL-001}} \newline \emph{Mandatory Optical Inter-Satellite Link Network} & The \abbr{SSEG} \textbf{shall} incorporate a mandatory Optical Inter-Satellite Link (\abbr{OISL}) network interconnecting the 72 \texttt{SATL} spacecraft in an intra-/inter-plane mesh and linking the \abbr{LEO} tracking layer to the 6 \texttt{SATG} geostationary relay nodes, guaranteeing tracking command delivery (\texttt{R-044}), inter-satellite \abbr{LOS} vector exchange (\texttt{R-023}), and observation downlink (\texttt{R-033}) regardless of spacecraft geographic position, in compliance with the link budgets in \textbf{RD-02} and subsystem specifications in \textbf{RD-03}/\textbf{RD-04}. & \texttt{R-023}, \texttt{R-033}, \texttt{R-044} & \texttt{T, A} \\ \hline
\textbf{\texttt{SSEG-REQ-ISL-002}} \newline \emph{All-Weather} \newline \emph{Space-to-Ground} \newline \emph{RF Links} & Each spacecraft \textbf{shall} provide redundant all-weather radio-frequency (\abbr{RF}) Telemetry, Tracking, and Command (\abbr{TT\&C}) and mission data links with the Ground Station Network (\texttt{123456.02.05}) supporting transmission of event reports, observation packages, and constellation control in accordance with \textbf{RD-03} and \textbf{RD-04}. & \texttt{R-030}, \texttt{R-033}, \texttt{R-042} & \texttt{T, A} \\ \hline
\textbf{\texttt{SSEG-REQ-ISL-003}} \newline \emph{End-to-End} \newline \emph{Cryptographic Security and Authentication} & All space-to-ground and inter-satellite communication links \textbf{shall} employ end-to-end encryption and telecommand authentication mechanisms compliant with applicable security requirements. & \texttt{R-003}, \texttt{R-031}, \texttt{R-034}, \texttt{R-035} & \texttt{T, R} \\ \hline
\textbf{\texttt{SSEG-REQ-ISL-004}} \newline \emph{Anti-Jamming Protection and Autonomous Mesh Failover} & \abbr{SSEG} communication links and \abbr{OISL} routing protocols \textbf{shall} incorporate protection against jamming, spoofing, and unauthorized access (\texttt{R-032}), and \textbf{shall} autonomously prioritize warning-chain traffic and reroute packets around failed nodes or degraded links. & \texttt{R-017}, \texttt{R-032}, \texttt{R-033}, \texttt{R-044} & \texttt{T, A} \\ \hline
\end{longtable}}

% -----------------------------------------------------------------------------
% 7.6 SPACE ENVIRONMENT, SURVIVABILITY, AND SUSTAINABILITY REQUIREMENTS
% -----------------------------------------------------------------------------
\subsubsection{7.6 Space Environment, Survivability, and Sustainability Requirements (\texttt{SSEG-REQ-ENV})}

Table~\ref{tab:req_env} specifies the segment-level orbital environment survivability, mandatory active end-of-life disposal, and passivation requirements.

{\footnotesize\setlength{\tabcolsep}{4pt}\renewcommand{\arraystretch}{1.15}%
\arrayrulecolor{TableBorderBlue}%
\rowcolors{2}{white}{TableRowLightBlue}%
\setlength{\LTcapwidth}{\textwidth}%
\begin{longtable}{|>{\raggedright\arraybackslash}p{0.18\textwidth}|>{\raggedright\arraybackslash}p{0.58\textwidth}|>{\raggedright\arraybackslash}p{0.12\textwidth}|>{\raggedright\arraybackslash}p{0.06\textwidth}|}
\caption{Space Environment, Survivability, and Sustainability Requirements (\texttt{SSEG-REQ-ENV})}\label{tab:req_env}\\
\hline
\rowcolor{TableHeaderBlue}
\thead{Req.\ ID \& Title} & \thead{Normative Requirement Statement (\texttt{shall})} & \thead{Parent Req.} & \thead{Verif.} \\ \hline
\endfirsthead
\rowcolor{white}\multicolumn{4}{l}{\footnotesize\textit{Table \thetable\ (continued): Space Environment, Survivability, and Sustainability Requirements (\texttt{SSEG-REQ-ENV})}} \\[0.3ex]
\hline
\rowcolor{TableHeaderBlue}
\thead{Req.\ ID \& Title} & \thead{Normative Requirement Statement (\texttt{shall})} & \thead{Parent Req.} & \thead{Verif.} \\ \hline
\endhead
\rowcolor{white}\multicolumn{4}{r}{\footnotesize\textit{Continued on next page}} \\
\endfoot
\endlastfoot
\textbf{\texttt{SSEG-REQ-ENV-001}} \newline \emph{Orbital Environment Survivability} & All \texttt{SATL} and \texttt{SATG} flight systems \textbf{shall} withstand the ionizing radiation, single-event effects, plasma charging, thermal/eclipse cycling, hard vacuum, and micrometeoroid/orbital debris (\abbr{MMOD}) environments of the 1,000~km \abbr{LEO} and 35,786~km \abbr{GEO} regimes over their respective operational lifetimes in accordance with the environmental budgets in \textbf{RD-02} and subsystem specifications in \textbf{RD-03}/\textbf{RD-04}. & \texttt{R-015}, \texttt{R-016}, \texttt{R-017} & \texttt{T, A} \\ \hline
\textbf{\texttt{SSEG-REQ-ENV-002}} \newline \emph{Mandatory Active End-of-Life Disposal} & In compliance with \texttt{R-018} and space debris mitigation guidelines, every \texttt{SATL} spacecraft \textbf{shall} execute active deorbitation from 1,000~km at End of Life to ensure complete atmospheric demise within $< 5\text{ years}$ of mission completion, and every \texttt{SATG} spacecraft \textbf{shall} execute end-of-life re-orbitation to a super-synchronous graveyard orbit at least $\Delta H \ge +300\text{ km}$ above 35,786~km. & \texttt{R-018} & \texttt{A, R} \\ \hline
\textbf{\texttt{SSEG-REQ-ENV-003}} \newline \emph{End-of-Life Constellation Passivation} & Upon completing end-of-life disposal manoeuvres, each \texttt{SATL} and \texttt{SATG} spacecraft \textbf{shall} permanently passivate all onboard stored energy sources, including venting/depleting residual propellants, disconnecting battery charging circuits, spinning down momentum actuators, and disabling all \abbr{RF} and optical transmitters. & \texttt{R-018} & \texttt{R, T} \\ \hline
\end{longtable}}

% -----------------------------------------------------------------------------
% 7.7 CONSTELLATION OPERATIONS, AUTONOMY, AND EXTERNAL INTERFACE REQUIREMENTS
% -----------------------------------------------------------------------------
\subsubsection{7.7 Constellation Operations, Autonomy, and External Interface Requirements (\texttt{SSEG-REQ-OPS} \& \texttt{SSEG-REQ-IF})}

Table~\ref{tab:req_ops_if} defines the constellation-level operational lifetime, redundancy, \abbr{FDIR} autonomy, event buffering, and external interface (\texttt{IF-EXT-01}--\texttt{IF-EXT-04}) requirements.

{\footnotesize\setlength{\tabcolsep}{4pt}\renewcommand{\arraystretch}{1.15}%
\arrayrulecolor{TableBorderBlue}%
\rowcolors{2}{white}{TableRowLightBlue}%
\setlength{\LTcapwidth}{\textwidth}%
\begin{longtable}{|>{\raggedright\arraybackslash}p{0.18\textwidth}|>{\raggedright\arraybackslash}p{0.58\textwidth}|>{\raggedright\arraybackslash}p{0.12\textwidth}|>{\raggedright\arraybackslash}p{0.06\textwidth}|}
\caption{Constellation Operations, Autonomy, and External Interface Requirements (\texttt{SSEG-REQ-OPS} \& \texttt{SSEG-REQ-IF})}\label{tab:req_ops_if}\\
\hline
\rowcolor{TableHeaderBlue}
\thead{Req.\ ID \& Title} & \thead{Normative Requirement Statement (\texttt{shall})} & \thead{Parent Req.} & \thead{Verif.} \\ \hline
\endfirsthead
\rowcolor{white}\multicolumn{4}{l}{\footnotesize\textit{Table \thetable\ (continued): Constellation Operations, Autonomy, and External Interface Requirements (\texttt{SSEG-REQ-OPS} \& \texttt{IF})}} \\[0.3ex]
\hline
\rowcolor{TableHeaderBlue}
\thead{Req.\ ID \& Title} & \thead{Normative Requirement Statement (\texttt{shall})} & \thead{Parent Req.} & \thead{Verif.} \\ \hline
\endhead
\rowcolor{white}\multicolumn{4}{r}{\footnotesize\textit{Continued on next page}} \\
\endfoot
\endlastfoot
\textbf{\texttt{SSEG-REQ-OPS-001}} \newline \emph{Operational Design Lifetimes} & The \abbr{SSEG} \textbf{shall} be designed and qualified for a minimum operational lifetime of 7~years for the \abbr{LEO} satellite constellation (\texttt{SATL}) and 15~years for the \abbr{GEO} satellite constellation (\texttt{SATG}). & \texttt{R-016} & \texttt{A, R} \\ \hline
\textbf{\texttt{SSEG-REQ-OPS-002}} \newline \emph{No Single-Point Failure Redundancy and Autonomous Safe Mode} & Critical spacecraft and constellation functional chains \textbf{shall} implement redundant architectures with no single-point failure (\texttt{R-017}) and autonomous Fault Detection, Isolation, and Recovery (\abbr{FDIR}) logic that switches to redundant hardware or transitions a severely faulted spacecraft into a power- and optically-safe mode without interrupting constellation surveillance. & \texttt{R-015}, \texttt{R-017}, \texttt{R-030} & \texttt{T, A} \\ \hline
\textbf{\texttt{SSEG-REQ-OPS-003}} \newline \emph{Health Telemetry, Onboard Event Buffering, and OTA Calibration} & Each spacecraft \textbf{shall} generate continuous health and status telemetry for Mission Control Center monitoring (\texttt{R-030}), maintain encrypted non-volatile onboard buffering of event observations and audit logs prior to downlink to guarantee zero data loss for 5-year archival (\texttt{R-004}, \texttt{R-035}, \texttt{R-036}, \texttt{R-047}), and support authenticated over-the-air (\abbr{OTA}) updates of onboard detection and tracking algorithms (\texttt{R-004}, \texttt{R-041}). & \texttt{R-004}, \texttt{R-030}, \texttt{R-035}, \texttt{R-036}, \texttt{R-041}, \texttt{R-047} & \texttt{T, R} \\ \hline
\textbf{\texttt{SSEG-REQ-IF-001}} \newline \emph{Space-to-Ground} \newline \emph{Interface Compliance (\texttt{IF-EXT-01})} & All \abbr{SSEG} housekeeping telemetry, tracking commands, \emph{GEO Event Reports}, \emph{LEO Observation Packages}, and \emph{Onboard Preliminary Trajectories} exchanged across \texttt{IF-EXT-01} \textbf{shall} conform to the packet schemas, timing, and cryptographic boundaries defined in the Space Segment \abbr{ICD} (\texttt{123456.05.02.04}) and Ground Segment Data Dictionary (Chapter~12 of \textbf{AD-01}). & \texttt{R-003}, \texttt{R-031}, \texttt{R-033}, \texttt{R-042}--\texttt{R-045} & \texttt{T, I} \\ \hline
\textbf{\texttt{SSEG-REQ-IF-002}} \newline \emph{Launch Vehicle and Dispenser Interface Compliance (\texttt{IF-EXT-02})} & All \texttt{SATL} and \texttt{SATG} mechanical, electrical, dynamic load, and launch safety-inhibit interfaces across \texttt{IF-EXT-02} \textbf{shall} comply with the multi-satellite plane-batch rideshare dispensers (\texttt{SATL}) and heavy-lift launch vehicle adapters (\texttt{SATG}) specified in \textbf{RD-01}, \textbf{RD-03}, and \textbf{RD-04}. & \texttt{R-005} & \texttt{I, T} \\ \hline
\textbf{\texttt{SSEG-REQ-IF-003}} \newline \emph{Target Plume Radiometric Interface Compliance (\texttt{IF-EXT-03})} & The \abbr{SSEG} optical payloads and onboard discrimination algorithms across \texttt{IF-EXT-03} \textbf{shall} detect and track boost- and post-boost-phase missile exhaust plumes against sunlit cloud tops, Earth-limb airglow, and terrestrial thermal clutter backgrounds. & \texttt{R-009}, \texttt{R-019}--\texttt{R-022}, \texttt{R-025} & \texttt{T, A} \\ \hline
\textbf{\texttt{SSEG-REQ-IF-004}} \newline \emph{Space Traffic Management and Regulatory Interface (\texttt{IF-EXT-04})} & The \abbr{SSEG} \textbf{shall} ingest Conjunction Data Messages (\abbr{CDM}s) relayed from the Ground Segment Mission Control Center across \texttt{IF-EXT-04} to execute collision-avoidance manoeuvres, and \textbf{shall} comply with \abbr{ITU} spectral regulations, laser safety interlocks, and space-debris disposal requirements. & \texttt{R-001}, \newline \texttt{R-018} & \texttt{T, R} \\ \hline
\end{longtable}}

% -----------------------------------------------------------------------------
% 7.8 SPACE SEGMENT VERIFICATION REQUIREMENTS
% -----------------------------------------------------------------------------
\subsubsection{7.8 Space Segment Verification Requirements (\texttt{SSEG-REQ-VER})}

Table~\ref{tab:req_ver} specifies the mandatory simulation, Hardware-in-the-Loop (\abbr{HIL}), environmental qualification, in-orbit validation, and archive-replay verification requirements for \abbr{SSEG} (\texttt{R-038}--\texttt{R-041}, \textbf{AD-02}). Detailed requirement-level verification execution specifications are provided in Section~7 of the Verification Plan (\texttt{EWAS-HSB-VP-01}), and requirement-by-requirement verification control status is recorded in the Verification Control Document (\texttt{EWAS-HSB-VCD-01}).

{\footnotesize\setlength{\tabcolsep}{4pt}\renewcommand{\arraystretch}{1.15}%
\arrayrulecolor{TableBorderBlue}%
\rowcolors{2}{white}{TableRowLightBlue}%
\setlength{\LTcapwidth}{\textwidth}%
\begin{longtable}{|>{\raggedright\arraybackslash}p{0.18\textwidth}|>{\raggedright\arraybackslash}p{0.58\textwidth}|>{\raggedright\arraybackslash}p{0.12\textwidth}|>{\raggedright\arraybackslash}p{0.06\textwidth}|}
\caption{Space Segment Verification Technical Requirements (\texttt{SSEG-REQ-VER})}\label{tab:req_ver}\\
\hline
\rowcolor{TableHeaderBlue}
\thead{Req.\ ID \& Title} & \thead{Normative Requirement Statement (\texttt{shall})} & \thead{Parent Req.} & \thead{Verif.} \\ \hline
\endfirsthead
\rowcolor{white}\multicolumn{4}{l}{\footnotesize\textit{Table \thetable\ (continued): Space Segment Verification Technical Requirements (\texttt{SSEG-REQ-VER})}} \\[0.3ex]
\hline
\rowcolor{TableHeaderBlue}
\thead{Req.\ ID \& Title} & \thead{Normative Requirement Statement (\texttt{shall})} & \thead{Parent Req.} & \thead{Verif.} \\ \hline
\endhead
\rowcolor{white}\multicolumn{4}{r}{\footnotesize\textit{Continued on next page}} \\
\endfoot
\endlastfoot
\textbf{\texttt{SSEG-REQ-VER-001}} \newline \emph{Multi-Scenario Constellation and Radiometric Simulation} & \abbr{SSEG} compliance with 24/7 Baltic and extra-regional coverage (\texttt{R-006}--\texttt{R-008}, \texttt{R-049}), boost-phase detection probability ($P_d > 95\%$, \texttt{R-010}), stereoscopic viewing geometry (\texttt{R-023}), and the $\le 60\text{ s}$ end-to-end warning latency budget (\texttt{R-011}) \textbf{shall} be verified via astrodynamic and radiometric simulation campaigns across representative ballistic missile launch trajectories and threat profiles in accordance with \textbf{AD-02}. & \texttt{R-038}, \newline \texttt{AD-02} & \texttt{A} \\ \hline
\textbf{\texttt{SSEG-REQ-VER-002}} \newline \emph{Hardware-in-the-Loop Testing of Sensing, LOS, and OBDH/OISL} & Prior to flight deployment, the \abbr{SWIR}/\abbr{MWIR} focal plane readout, onboard proton-despiking and stereoscopic 3D triangulation pipeline (\texttt{R-022}, \texttt{R-023}), independent payload optical \abbr{LOS} $< 5\text{ s}$ lock-on and $< 2\text{ arcsec/s}$ stabilization loop (\texttt{R-024}), and encrypted \abbr{OISL} routing protocols (\texttt{R-031}, \texttt{R-044}) \textbf{shall} undergo closed-loop Hardware-in-the-Loop (\abbr{HIL}) testing on the Engineering and Qualification Models (\abbr{WBS} \texttt{123456.04.02.02}). & \texttt{R-039}, \newline \texttt{AD-02} & \texttt{T} \\ \hline
\textbf{\texttt{SSEG-REQ-VER-003}} \newline \emph{Environmental Qualification and Batch Acceptance Testing} & The lead \texttt{SATL} and \texttt{SATG} Qualification / Proto-Flight Models (\texttt{QM}/\texttt{PFM}) and all recurring Flight Models (\texttt{FM-01} to \texttt{FM-72} for \texttt{SATL}; \texttt{FM-01} to \texttt{FM-06} for \texttt{SATG}) \textbf{shall} undergo mechanical vibration, shock, acoustic, thermal-vacuum (\abbr{TVAC}), and electromagnetic compatibility (\abbr{EMC}) testing in accordance with the project environmental qualification and acceptance rules (\textbf{AD-02}). & \texttt{R-015}, \texttt{R-016}, \texttt{AD-02} & \texttt{T} \\ \hline
\textbf{\texttt{SSEG-REQ-VER-004}} \newline \emph{End-to-End} \newline \emph{Operational Validation} & During In-Orbit Commissioning and periodic operational readiness exercises, \abbr{SSEG} and the Ground Segment \textbf{shall} execute end-to-end validation trials using cooperative calibration targets and injected or live rocket-launch signatures to demonstrate compliance with the $\le 60\text{ s}$ total warning latency (\texttt{R-011}), $P_d > 95\%$ (\texttt{R-010}), and $> 99.9\%$ availability (\texttt{R-015}). & \texttt{R-040}, \newline \texttt{AD-02} & \texttt{T} \\ \hline
\textbf{\texttt{SSEG-REQ-VER-005}} \newline \emph{Post-Event Archive Analysis and Algorithm Calibration} & The \abbr{SSEG} flight software and Ground Segment support tools \textbf{shall} support periodic ingestion and replay of archived mission event data (under \texttt{R-004}) to verify false-alarm reduction and update onboard plume-classification coefficients via authenticated over-the-air patches. & \texttt{R-004}, \newline \texttt{R-041} & \texttt{A, T} \\ \hline
\end{longtable}}
